\documentclass{article}
\usepackage{amsmath,amssymb,amsthm}
\newtheorem{theorem}{Theorem}
\newtheorem{lemma}[theorem]{Lemma}
\newtheorem{corollary}[theorem]{Corollary}
\newcounter{paperpremise}
\newcounter{definingequation}
\newenvironment{proofstep}{\par\smallskip\ifdim\leftskip<3em\relax\advance\leftskip by 1em\relax\fi\noindent\ignorespaces}{\par\smallskip}
\setlength{\textwidth}{6.5in}
\setlength{\oddsidemargin}{0in}
\setlength{\emergencystretch}{\linewidth}
\title{Counting 580}
\author{}
\date{}
\begin{document}
\maketitle
\section*{Definitions}
\paragraph{Counts.} \par\noindent\refstepcounter{definingequation}\label{defeq:1}\textup{(D1)} \(\operatorname{counts}\left(0\right) = \left(1, 1\right)\).
\par\noindent\refstepcounter{definingequation}\label{defeq:2}\textup{(D2)} Let \(n\) be a natural number. Then \(\operatorname{counts}(n + 1) = (\pi_1(\operatorname{counts}(n)) + \pi_2(\operatorname{counts}(n)),\allowbreak{} \pi_1(\operatorname{counts}(n)))\).

\paragraph{Open count (\(A\)).} \par\noindent\refstepcounter{definingequation}\label{defeq:3}\textup{(D3)} Let \(n\) be a natural number. Then \(A\left(n\right) = \pi_1(\operatorname{counts}\left(n\right))\).
\par\noindent\refstepcounter{definingequation}\label{defeq:4}\textup{(D4)} \(A\left(0\right) = 1\).

\paragraph{Forced count (\(B\)).} \par\noindent\refstepcounter{definingequation}\label{defeq:5}\textup{(D5)} Let \(n\) be a natural number. Then \(B\left(n\right) = \pi_2(\operatorname{counts}\left(n\right))\).
\par\noindent\refstepcounter{definingequation}\label{defeq:6}\textup{(D6)} \(B\left(0\right) = 1\).

\subsection*{Prerequisites}
\begin{description}
\item[P1.] \refstepcounter{paperpremise}\label{premise:4e61742e6c655f7265666c}\textbf{Reflexivity of order.} Let \(n\) be a natural number. Then \(n \le n\).
\item[P2.] \refstepcounter{paperpremise}\label{premise:4e61742e6c655f6164645f7269676874}\textbf{Adding a natural number can only increase a sum.} Let \(n\) and \(k\) be natural numbers. Then \(n \le n + k\).
\item[P3.] \refstepcounter{paperpremise}\label{premise:4e61742e6164645f6c655f6164645f7269676874}\textbf{Adding a common right summand preserves order.} Let \(n\), \(m\) and \(k\) be natural numbers. If \(n \le m\), then \(n + k \le m + k\).
\item[P4.] \refstepcounter{paperpremise}\label{premise:4e61742e6164645f6c655f6164645f6c656674}\textbf{Adding a common left summand preserves order.} Let \(n\), \(m\) and \(k\) be natural numbers. If \(n \le m\), then \(k + n \le k + m\).
\item[P5.] \refstepcounter{paperpremise}\label{premise:4e61742e6d756c5f74776f}\textbf{Doubling by multiplication.} Let \(n\) be a natural number. Then \(n \cdot 2 = n + n\).
\item[P6.] \refstepcounter{paperpremise}\label{premise:4e61742e6d756c5f6c655f6d756c5f7269676874}\textbf{Multiplying by a common natural number preserves order.} Let \(n\), \(m\) and \(k\) be natural numbers. If \(n \le m\), then \(n \cdot k \le m \cdot k\).
\end{description}

\section{Theorems}
\begin{lemma}[{Split at first bit}]\label{theorem:1}
Let \(n\) be a natural number. Then \(A(n + 1) = A(n) + B(n)\).
\end{lemma}

\begin{proof}
Evaluating the two sides gives
\[
A(n + 1) = A(n) + B(n).
\]
\end{proof}

\begin{lemma}[{Forced zero}]\label{theorem:2}
Let \(n\) be a natural number. Then \(B(n + 1) = A(n)\).
\end{lemma}

\begin{proof}
Evaluating the two sides gives
\[
B(n + 1) = A(n).
\]
\end{proof}

\begin{corollary}[{Two step recurrence}]\label{theorem:3}
Let \(n\) be a natural number. Then \(A(n + 2) = A(n + 1) + A(n)\).
\end{corollary}

\begin{proof}
\begin{align*}
A(n + 2) &= A(n + 1) + B(n + 1) && \text{by Lemma~\ref{theorem:1}} \\
 &= A(n + 1) + A(n). && \text{by Lemma~\ref{theorem:2}}
\end{align*}
\end{proof}

\begin{lemma}[{Forced at most open}]\label{theorem:4}
Let \(n\) be a natural number. Then \(B(n) \le A(n)\).
\end{lemma}

\begin{proof}
We distinguish the zero and successor cases for \(n\).
\begin{description}
\item[Zero case.] Applying \emph{Reflexivity of order} gives \(1 \le 1\).
By the defining equations, this gives \(B(0) \le A(0)\).
\item[Successor case.] Let \(i\) be a natural number.
\begin{align*}
B(i + 1) &= A(i) && \text{by Lemma~\ref{theorem:2}} \\
 &\le A(i) + B(i) && \text{by Prerequisite~P\ref{premise:4e61742e6c655f6164645f7269676874}} \\
 &= A(i + 1). && \text{by Lemma~\ref{theorem:1}}
\end{align*}
\end{description}
\end{proof}

\begin{lemma}[{Two bit bound}]\label{theorem:5}
Let \(n\) be a natural number. Then \(A(n + 2) \le A(n) \cdot 3\).
\end{lemma}

\begin{proof}
Lemma~\ref{theorem:4} gives \(B(n) \le A(n)\).
Apply Corollary~\ref{theorem:3} and Lemma~\ref{theorem:1} in order, use monotonicity of addition and collect equal summands.
\[
\begin{aligned}
A(n + 2) &= A(n + 1) + A(n) \\
 &= A(n) + B(n) + A(n) \\
 &\le A(n) + A(n) + A(n) \\
 &= A(n) \cdot 3.
\end{aligned}
\]
\end{proof}

\begin{theorem}[{Exponential bound}]\label{theorem:6}
Let \(k\) be a natural number. Then \(A(2 \cdot k) \le 3^{k}\).
\end{theorem}

\begin{proof}
We proceed by induction on \(k\).
\begin{description}
\item[Base case.] Applying \emph{Reflexivity of order} gives \(1 \le 1\).
By the defining equations, this gives
\[
A(2 \cdot 0) \le 3^{0}.
\]
\item[Inductive step.] Let \(i\) be a natural number.
As the induction hypothesis, assume
\[
A(2 \cdot i) \le 3^{i}.
\]
\begin{align*}
A(2 \cdot (i + 1)) &\le A(2 \cdot i) \cdot 3 && \text{by Lemma~\ref{theorem:5}} \\
 &\le 3^{i + 1}. && \text{by Prerequisite~P\ref{premise:4e61742e6d756c5f6c655f6d756c5f7269676874} applied to the induction hypothesis}
\end{align*}
\end{description}
This completes the induction.
\end{proof}

\end{document}
